\begin{table}[htbp]
\centering
\small
\caption{Club-price rights ratios $\rho_i$ equivalent to the Wolfram et al. and Banerjee, Duflo \& Greenstone proposals --- option B (simultaneous solve)}
\renewcommand{\arraystretch}{1.15}
\begin{tabular}{lcccc}
  \toprule
  & \multicolumn{2}{c}{\textbf{Wolfram et al.}} & \multicolumn{2}{c}{\textbf{Banerjee, Duflo \& Greenstone}} \\
  \cmidrule(lr){2-3} \cmidrule(lr){4-5}
  \textbf{Country} & $p_i$ \textbf{2030 (\$/t)} & $\rho_i$ & $p_i$ \textbf{2030 (\$/t)} & $\rho_i$ \\
  \midrule
  United States & 0 & -- & 75 & 2.833 \\
  Russia & 0 & -- & 75 & 1.786 \\
  European Union & 75 & 0.475 & 75 & 0.454 \\
  Turkey & 0 & -- & 50 & 1.128 \\
  China & 50 & 0.812 & 50 & 0.820 \\
  Indonesia & 50 & 0.716 & 50 & 0.774 \\
  Brazil & 50 & 0.268 & 50 & 0.291 \\
  India & 25 & 0.594 & 30 & 0.622 \\
  Nigeria & 0 & -- & 30 & 0.123 \\
  DR Congo & 0 & -- & 10 & 0.026 \\
  \midrule
  \multicolumn{5}{l}{\textit{Variant 1 --- unscaled rights $\rho_i$ (the allocation sets the cap)}} \\
  \quad Emissions change in the coalition (\%) &  & $-$34.5\% &  & $-$27.2\% \\
  \quad World temperature 2100, change ($^\circ$C) &  & $-$0.068 &  & $-$0.108 \\
  \quad World welfare gain, EDE (\%) &  & 0.322\% &  & 0.618\% \\
  \quad World mean consumption gain (\%) &  & $-$0.012\% &  & $-$0.020\% \\
  \midrule
  \multicolumn{5}{l}{\textit{Variant 2 --- surplus allocated to minimise the largest shortfall}} \\
  \quad World welfare gain, EDE (\%) &  & 0.025\% &  & 0.224\% \\
  \quad World mean consumption gain (\%) &  & 0.016\% &  & 0.045\% \\
  \midrule
  \multicolumn{5}{l}{\textit{Variant 3 --- surplus shared by marginal utility}} \\
  \quad World welfare gain, EDE (\%) &  & 0.650\% &  & 1.323\% \\
  \quad World mean consumption gain (\%) &  & 0.031\% &  & 0.055\% \\
  \midrule
  \multicolumn{5}{l}{\textit{Variant 4 --- surplus shared by uniform scaling}} \\
  \quad World welfare gain, EDE (\%) &  & $-$0.057\% &  & 0.063\% \\
  \quad World mean consumption gain (\%) &  & 0.013\% &  & 0.040\% \\
  \bottomrule
\end{tabular}
\label{tab:equiv_rights_B}
\\[4pt]
{\footnotesize Note: $p_i$ is the carbon price the proposal asks of country $i$ in 2030; the comparison regime instead prices every member of the proposal's own club at one common price (\$49.2/t for Wolfram et al. and \$55.8/t for Banerjee, Duflo \& Greenstone in 2030) and leaves non-members unpriced, as the proposal does. $\rho_i$ is the country's allocation as a multiple of an equal-per-capita share of the club's emissions under the proposal; the European Union figure aggregates its members' rights. A country the proposal does not price is outside the club in both regimes: it is shown at $p_i=0$ with no equivalent allocation (`--'). $\rho_i < 0$ would mean a country prefers the club regime even while paying a net fee for its allocation. Variant 1 lets the equivalent allocation set the cap, so the club price is recalibrated to deliver exactly those rights; variant 2 rescales the same allocation to the club's own emissions under the proposal, so club emissions are held fixed and only allocative efficiency remains; it is reported both on the inequality-averse global EDE ($\eta=1.5$) and on world mean consumption per capita, which weights every person's consumption equally. The temperature row is the change from the proposal's own 2100 warming, which is 2.16\textdegree{}C for Wolfram et al. and 1.80\textdegree{}C for Banerjee, Duflo \& Greenstone. Both proposals are priced from 2025 and capped at the backstop price. The Banerjee et al. proposal prices high-income countries at \$75/t, a tier the report leaves unspecified; the legacy column, where it appears, leaves them unpriced.}
\end{table}
